% Owner: Hoang.
\section{Experimental Evaluation Plan}\label{sec:evaluation}

This section defines a prospective evaluation of the specified acceptance-to-draft workflow. The intended sequence is local acceptance plus an outbox row in Listing and Booking, publication of \texttt{RentalRequestAccepted} to \texttt{listing.events}, and draft creation by Contract and Tenant. The inspected backend does not yet provide the complete publisher and Compose topology for this sequence. The protocol must therefore be run only after an implementation revision and deployment manifest are fixed and recorded. No latency, throughput, recovery, or end-to-end consistency measurements are reported here.

\subsection{Experimental Setup}
Use a dedicated deployment that includes PostgreSQL, Kafka, the gateway, Listing and Booking, and Contract and Tenant; the current Compose file must first be extended or replaced. Record the source revision, deployment manifest, image versions, broker and database settings, CPU and memory limits, network placement, and load-generator location. Prepare synthetic accounts with the required roles, active listings, and eligible requests. Use distinct request identifiers and restore equivalent fixture data for each independent repetition. Before inspecting reported runs, freeze the offered rates, operation mix, warm-up, measurement duration, timeout, event observation cutoff, and fault triggers. Run at least three independent repetitions as a minimum feasibility check; choose a larger number and batch size when needed to estimate tail-latency or rare-failure uncertainty. Report every run, including failures, and state what precision the sample supports.

\subsection{Boundary and Authorization Audit (RQ1)}
For each of the seven proposed services, record its authoritative tables, configured database user, inbound APIs, outgoing calls, and published and consumed event types at the tested revision. Search source and configuration for cross-database connections and verify grants using service credentials against a fresh database volume. The boundary criterion is no service write to another service's tables in the audited paths; an untested path is reported as unverified, not compliant. Test the acceptance path and an authorization matrix for CUSTOMER, SALE, QUAN\_LY, KY\_THUAT, and ADMIN against owned and non-owned records and assigned and unassigned buildings. Include direct access to the loopback-published User service and forged identity headers as negative tests. The audit gives RQ1 a checkable answer; it does not prove that the chosen partition is optimal.

\subsection{API Latency and Throughput}
Test public listing reads and authenticated rental-request writes through the gateway after both routes are deployed. Measure each operation separately and keep the operation mix fixed for a combined workload. A pilot should identify feasible offered rates, including light load and the onset of sustained queueing or failures; record the rule used to choose the final rates. Freeze rate steps, concurrency, warm-up, measurement duration, query distribution, and request timeout before reported runs. Report offered and successful requests per second, counts of successful responses, HTTP errors, and timeouts, and p50, p95, and p99 of successful response latency with sample counts and run-to-run variation. Timeouts remain in the failure denominator. These measurements characterize the tested deployment; an architectural cost comparison requires a separately implemented, resource-matched baseline.

\subsection{Event Delivery and Duplicate Handling}
Accept a known batch of verified requests. Use each committed acceptance as the denominator and reconcile it with its outbox row, all matching broker records, and all draft contracts by rental request and event identifiers. Record a correlation identifier when present. Capture acceptance-commit, publication, and effect-commit times; document any timestamp proxy and its resolution. Report commit-to-publication and commit-to-contract latency separately. At a cutoff fixed before the run, classify each acceptance as: completed (at least one matching publication and exactly one correct draft), pending relay (outbox row not published), published but awaiting a consumer effect, dead-lettered, or unexplained. A request with multiple drafts or incorrect business fields is a failed invariant even if one matching draft exists. Report counts and identifiers in every class; do not treat pending or dead-lettered records as successful recovery.

Replay the original Kafka envelope with the same \texttt{eventId} and key after verifying that the first delivery committed a draft. Repeat after restarting Contract and Tenant. Verify receipt, the processed-event row, one draft per request, and unchanged draft fields. Inspect any other effect attributed to this event in the tested consumer; notification delivery is a separate consumer outcome and must be assessed separately if included in a broader claim. Report replay attempts, observed broker deliveries, and duplicate business effects with distinct denominators. Also test whether the retained synchronous \texttt{/internal/contracts/drafts} path can create a second draft for the same rental request.

\subsection{Fault Tolerance and Recovery}
Run the same workflow without injected faults to establish a reference completion distribution. For a consumer outage, stop Contract and Tenant, accept a batch, confirm broker publication, restore the consumer, and observe the backlog. For a publisher outage, interrupt the relay after acceptance commits and before publication; verify pending outbox rows, then restore it. A separate failpoint after broker acknowledgement but before \texttt{published\_at} is recorded is required to evaluate duplicate publication. If that failpoint cannot be implemented reliably, label the scenario untested and make no claim about that crash window. Fault injection should confirm that the intended component and interval were reached \cite{meiklejohn2021faultinjection}.

For each fault, record injection, restoration, component readiness, and the last matching draft commit. Report both injection-to-completion time, which includes deliberate downtime, and restoration-to-backlog-clearance time. At the predefined cutoff, the tested acceptance invariant passes only when every committed acceptance has at least one corresponding published event, exactly one correct draft, and no unresolved dead-letter record. Otherwise report the pending, dead-lettered, duplicate-effect, and unexplained cases separately. A finite trial cannot establish eventual completion beyond its cutoff. Table~\ref{tab:experiments} summarizes the planned scenarios and evidence.

% Owner: Hoang. All outcomes and measurements remain TODO until experiments are run.
\begin{table*}[t]
\caption{Experimental Scenarios and Required Evidence}
\label{tab:experiments}
\centering\footnotesize
\begin{tabularx}{\textwidth}{@{}p{0.22\textwidth}X X@{}}
\toprule
Scenario & Procedure and controlled variable & Evidence to report \\
\midrule
API throughput and latency & Increase offered request rate with a fixed operation mix; repeat each load level & Achieved throughput, error rate, latency percentiles, run variation: TODO \\
Kafka event processing & Commit a known sequence of changes; track corresponding publication and consumption & End-to-end lag, missing events, retries, trace IDs: TODO \\
Duplicate event handling & Replay known identifiers before and after consumer restart & Number of business effects per identifier and deduplication records: TODO \\
Service failure recovery & Stop a selected service during a defined workflow; restore it & Failed requests, pending work, invariant recovery time: TODO \\
Transactional outbox recovery & Interrupt relay around commit, send, and acknowledgement points & Outbox state, repeated publications, eventual consumer effects: TODO \\
\bottomrule
\end{tabularx}
\end{table*}

